\subsubsubsection{Past Examples of Periodic Boundary Conditions}
\label{SEC:BC:Numerical:Periodic:Examples}

These are examples taken from GFR input files that have been used in the past to
run simulations for problems including periodic boundaries.
%
Note that the array indices used for \inptext{BC\_Input} are simply copied over
from the input files, meaning each index is relative to the other boundary
conditions that also existed in the input file it was taken from.
%
Part of the intent behind keeping these array indices as they originally were in
their respective input file is to demonstrate that the previously used variable
indices of $i$ and $j$ are can represent any array element of
\inptext{BC\_Input} provided that they link to each other when they form a pair
of periodic boundaries.
%
\begin{itemize}
  \item Streamwise and spanwise periodic channel flow
   %\begin{nmldescitem}[escapechar=@]
   %\begin{nmldescitem}[fontupper=\ttfamily\footnotesize,
    \begin{nmldescitem}[fontupper=\ttfamily\scriptsize,
                        listing options={escapeinside={!@}{!@}}]
&nml_bc
  BC_Input(1)%Name = "INFLOW"
  BC_Input(1)%Set_BC_Value = 81
  BC_Input(1)%Periodic%Matching_BC_Name = "OUTFLOW"
  BC_Input(1)%Periodic%Matching_BC_Idx = 2
  BC_Input(1)%Periodic%Translation(1) = 0.036576 !@\vspace{5pt}!@
  BC_Input(2)%Name = "OUTFLOW"
  BC_Input(2)%Set_BC_Value = 81
  BC_Input(2)%Periodic%Matching_BC_Name = "INFLOW"
  BC_Input(2)%Periodic%Matching_BC_Idx = 1
  BC_Input(2)%Periodic%Translation(1) = -0.036576 !@\vspace{5pt}!@
  BC_Input(5)%Name = "ZMIN"
  BC_Input(5)%Set_BC_Value = 81
  BC_Input(5)%Periodic%Matching_BC_Name = "ZMAX"
  BC_Input(5)%Periodic%Matching_BC_Idx = 6
  BC_Input(5)%Periodic%Translation(3) = 0.018288 !@\vspace{5pt}!@
  BC_Input(6)%Name = "ZMAX"
  BC_Input(6)%Set_BC_Value = 81
  BC_Input(6)%Periodic%Matching_BC_Name = "ZMIN"
  BC_Input(6)%Periodic%Matching_BC_Idx = 5
  BC_Input(6)%Periodic%Translation(3) = -0.018288
/
    \end{nmldescitem}
   %\FINALFIX{\newpage}
  \item Single rotor blade from spacecraft cabin ventilation
    fan\cite{Spiegel_2025}
    \begin{itemize}[before=\small]
      \item Rotor consists of nine blades spinning about the $x$-axis resulting
        in a rotationally periodic domain that is rotated 40\textdegree{} about
        the $x$-axis.
      \item This grid has two pairs of rotationally periodic boundaries: one for
        the primary computational domain, and a second for a sponge region
        appended onto the regular outflow boundary to damp out non-physical
        pressure waves from reflecting off the outflow boundary.
        The periodic boundary is split in this way for convenience, but could
        be combined.
    \end{itemize}
   %\begin{nmldescitem}[escapechar=@]
   %\begin{nmldescitem}[fontupper=\ttfamily\footnotesize,
    \begin{nmldescitem}[fontupper=\ttfamily\scriptsize,
                        listing options={escapeinside={!@}{!@}}]
&nml_bc
  BC_Input(8)%Name = "ThetaMin"
  BC_Input(8)%Set_BC_Value = 81
  BC_Input(8)%Periodic%Matching_BC_Name = "ThetaMax"
  BC_Input(8)%Periodic%Matching_BC_Idx = 9
  BC_Input(8)%Periodic%RotationAngle(1) = 40.0 !@\vspace{5pt}!@
  BC_Input(9)%Name = "ThetaMax"
  BC_Input(9)%Set_BC_Value = 81
  BC_Input(9)%Periodic%Matching_BC_Name = "ThetaMin"
  BC_Input(9)%Periodic%Matching_BC_Idx = 8
  BC_Input(9)%Periodic%RotationAngle(1) = -40.0 !@\vspace{5pt}!@
  BC_Input(10)%Name = "ThetaMinSponge"
  BC_Input(10)%Set_BC_Value = 81
  BC_Input(10)%Periodic%Matching_BC_Name = "ThetaMaxSponge"
  BC_Input(10)%Periodic%Matching_BC_Idx = 11
  BC_Input(10)%Periodic%RotationAngle(1) = 40.0 !@\vspace{5pt}!@
  BC_Input(11)%Name = "ThetaMaxSponge"
  BC_Input(11)%Set_BC_Value = 81
  BC_Input(11)%Periodic%Matching_BC_Name = "ThetaMinSponge"
  BC_Input(11)%Periodic%Matching_BC_Idx = 10
  BC_Input(11)%Periodic%RotationAngle(1) = -40.0
/
    \end{nmldescitem}
  \item Ceramic Matrix Composite 9\% thickness (CMC9) blade
    cascade\cite{Giel_2020} with spanwise and vertical
    periodicity\cite{Miki_2022}
   %\begin{nmldescitem}[escapechar=@]
   %\begin{nmldescitem}[fontupper=\ttfamily\footnotesize,
    \begin{nmldescitem}[fontupper=\ttfamily\scriptsize,
                        listing options={escapeinside={!@}{!@}}]
&nml_bc
  BC_Input(3)%Name = "periodic_yp13"
  BC_Input(3)%Set_BC_Value = 81
  BC_Input(3)%Periodic%Matching_BC_Name = "periodic_ym"
  BC_Input(3)%Periodic%Matching_BC_Idx = 4
  BC_Input(3)%Periodic%Translation(2) = -0.13 !@\vspace{5pt}!@
  BC_Input(4)%Name = "periodic_ym"
  BC_Input(4)%Set_BC_Value = 81
  BC_Input(4)%Periodic%Matching_BC_Name = "periodic_yp13"
  BC_Input(4)%Periodic%Matching_BC_Idx = 3
  BC_Input(4)%Periodic%Translation(2) = 0.13 !@\vspace{5pt}!@
  BC_Input(5)%Name = "periodic_zp013"
  BC_Input(5)%Set_BC_Value = 81
  BC_Input(5)%Periodic%Matching_BC_Name = "periodic_z0"
  BC_Input(5)%Periodic%Matching_BC_Idx = 6
  BC_Input(5)%Periodic%Translation(3) = -0.052 !@\vspace{5pt}!@
  BC_Input(6)%Name = "periodic_z0"
  BC_Input(6)%Set_BC_Value = 81
  BC_Input(6)%Periodic%Matching_BC_Name = "periodic_zp013"
  BC_Input(6)%Periodic%Matching_BC_Idx = 5
  BC_Input(6)%Periodic%Translation(3) = 0.052
/
    \end{nmldescitem}
  \item Taylor-Couette flow\cite{Wang_2021} with streamwise periodicity
   %\begin{nmldescitem}[escapechar=@]
   %\begin{nmldescitem}[fontupper=\ttfamily\footnotesize,
    \begin{nmldescitem}[fontupper=\ttfamily\scriptsize,
                        listing options={escapeinside={!@}{!@}}]
&nml_bc
  BC_Input(3)%Name = "Inflow"
  BC_Input(3)%Set_BC_Value = 81
  BC_Input(3)%Periodic%Matching_BC_Name = "Outflow"
  BC_Input(3)%Periodic%Matching_BC_Idx = 4
  BC_Input(3)%Periodic%Translation(1) = 6.283185307 !@\vspace{5pt}!@
  BC_Input(4)%Name = "Outflow"
  BC_Input(4)%Set_BC_Value = 81
  BC_Input(4)%Periodic%Matching_BC_Name = "Inflow"
  BC_Input(4)%Periodic%Matching_BC_Idx = 3
  BC_Input(4)%Periodic%Translation(1) = -6.283185307
/
    \end{nmldescitem}
\end{itemize}

%###############################################################################

%\begin{itemize}
%  \item Streamwise and spanwise periodic channel flow
%    \begin{InputExample}[before=\footnotesize]
%      %
%      \item[{\BCinptext[\texttt{1}]{Name = "INFLOW"}}]
%      \item[{\BCinptext[\texttt{1}]{Set\_BC\_Value = 81}}]
%      \item[{\BCinptext[\texttt{1}]{Periodic\%Matching\_BC\_Name = "OUTFLOW"}}]
%      \item[{\BCinptext[\texttt{1}]{Periodic\%Matching\_BC\_Idx = 2}}]
%      \item[{\BCinptext[\texttt{1}]{Periodic\%Translation(1) = 0.036576}}]
%      %
%      \item[{\BCinptext[\texttt{2}]{Name = "OUTFLOW"}}]\vspace{5pt}
%      \item[{\BCinptext[\texttt{2}]{Set\_BC\_Value = 81}}]
%      \item[{\BCinptext[\texttt{2}]{Periodic\%Matching\_BC\_Name = "INFLOW"}}]
%      \item[{\BCinptext[\texttt{2}]{Periodic\%Matching\_BC\_Idx = 1}}]
%      \item[{\BCinptext[\texttt{2}]{Periodic\%Translation(1) = -0.036576}}]
%      %
%      \item[{\BCinptext[\texttt{5}]{Name = "ZMIN"}}]\vspace{5pt}
%      \item[{\BCinptext[\texttt{5}]{Set\_BC\_Value = 81}}]
%      \item[{\BCinptext[\texttt{5}]{Periodic\%Matching\_BC\_Name = "ZMAX"}}]
%      \item[{\BCinptext[\texttt{5}]{Periodic\%Matching\_BC\_Idx = 6}}]
%      \item[{\BCinptext[\texttt{5}]{Periodic\%Translation(3) = 0.018288}}]
%      %
%      \item[{\BCinptext[\texttt{6}]{Name = "ZMAX"}}]\vspace{5pt}
%      \item[{\BCinptext[\texttt{6}]{Set\_BC\_Value = 81}}]
%      \item[{\BCinptext[\texttt{6}]{Periodic\%Matching\_BC\_Name = "ZMIN"}}]
%      \item[{\BCinptext[\texttt{6}]{Periodic\%Matching\_BC\_Idx = 5}}]
%      \item[{\BCinptext[\texttt{6}]{Periodic\%Translation(3) = -0.018288}}]
%      %
%    \end{InputExample}
%  \item Single rotor blade from spacecraft cabin ventilation
%    fan\cite{Spiegel_2025}
%    \begin{itemize}
%      \item Rotor consists of nine blades spinning about the $x$-axis resulting
%        in a rotationally periodic domain that is rotated 40\textdegree{} about
%        the $x$-axis.
%      \item This grid has two pairs of rotationally periodic boundaries: one for
%        the primary computational domain, and a second for a sponge region
%        appended onto the regular outflow boundary to damp out non-physical
%        pressure waves from reflecting off the outflow boundary.
%    \end{itemize}
%    \begin{InputExample}[before=\footnotesize]
%      %
%      \item[{\BCinptext[\texttt{8}]{Name = "ThetaMin"}}]
%      \item[{\BCinptext[\texttt{8}]{Set\_BC\_Value = 81}}]
%      \item[{\BCinptext[\texttt{8}]{Periodic\%Matching\_BC\_Name = "ThetaMax"}}]
%      \item[{\BCinptext[\texttt{8}]{Periodic\%Matching\_BC\_Idx = 9}}]
%      \item[{\BCinptext[\texttt{8}]{Periodic\%RotationAngle(1) = 40.0}}]
%      %
%      \item[{\BCinptext[\texttt{9}]{Name = "ThetaMax"}}]\vspace{5pt}
%      \item[{\BCinptext[\texttt{9}]{Set\_BC\_Value = 81}}]
%      \item[{\BCinptext[\texttt{9}]{Periodic\%Matching\_BC\_Name = "ThetaMin"}}]
%      \item[{\BCinptext[\texttt{9}]{Periodic\%Matching\_BC\_Idx = 8}}]
%      \item[{\BCinptext[\texttt{9}]{Periodic\%RotationAngle(1) = -40.0}}]
%      %
%      \item[{\BCinptext[\texttt{10}]{Name = "ThetaMinSponge"}}]\vspace{5pt}
%      \item[{\BCinptext[\texttt{10}]{Set\_BC\_Value = 81}}]
%      \item[{\BCinptext[\texttt{10}]{Periodic\%Matching\_BC\_Name = "ThetaMaxSponge"}}]
%      \item[{\BCinptext[\texttt{10}]{Periodic\%Matching\_BC\_Idx = 11}}]
%      \item[{\BCinptext[\texttt{10}]{Periodic\%RotationAngle(1) = 40.0}}]
%      %
%      \item[{\BCinptext[\texttt{11}]{Name = "ThetaMaxSponge"}}]\vspace{5pt}
%      \item[{\BCinptext[\texttt{11}]{Set\_BC\_Value = 81}}]
%      \item[{\BCinptext[\texttt{11}]{Periodic\%Matching\_BC\_Name = "ThetaMinSponge"}}]
%      \item[{\BCinptext[\texttt{11}]{Periodic\%Matching\_BC\_Idx = 10}}]
%      \item[{\BCinptext[\texttt{11}]{Periodic\%RotationAngle(1) = -40.0}}]
%      %
%    \end{InputExample}
%  \item Ceramic Matrix Composite 9\% thickness (CMC9) blade\cite{Giel_2020} with
%    spanwise and vertical periodicity\cite{Miki_2022}
%    \begin{InputExample}[before=\footnotesize]
%      %
%      \item[{\BCinptext[\texttt{3}]{Name = "periodic\_yp13"}}]
%      \item[{\BCinptext[\texttt{3}]{Set\_BC\_Value = 81}}]
%      \item[{\BCinptext[\texttt{3}]{Periodic\%Matching\_BC\_Name = "periodic\_ym"}}]
%      \item[{\BCinptext[\texttt{3}]{Periodic\%Matching\_BC\_Idx = 4}}]
%      \item[{\BCinptext[\texttt{3}]{Periodic\%Translation(2) = -0.13}}]
%      %
%      \item[{\BCinptext[\texttt{4}]{Name = "periodic\_ym"}}]\vspace{5pt}
%      \item[{\BCinptext[\texttt{4}]{Set\_BC\_Value = 81}}]
%      \item[{\BCinptext[\texttt{4}]{Periodic\%Matching\_BC\_Name = "periodic\_yp13"}}]
%      \item[{\BCinptext[\texttt{4}]{Periodic\%Matching\_BC\_Idx = 3}}]
%      \item[{\BCinptext[\texttt{4}]{Periodic\%Translation(2) = 0.13}}]
%      %
%      \item[{\BCinptext[\texttt{5}]{Name = "periodic\_zp013"}}]\vspace{5pt}
%      \item[{\BCinptext[\texttt{5}]{Set\_BC\_Value = 81}}]
%      \item[{\BCinptext[\texttt{5}]{Periodic\%Matching\_BC\_Name = "periodic\_z0"}}]
%      \item[{\BCinptext[\texttt{5}]{Periodic\%Matching\_BC\_Idx = 6}}]
%      \item[{\BCinptext[\texttt{5}]{Periodic\%Translation(3) = -0.052}}]
%      %
%      \item[{\BCinptext[\texttt{6}]{Name = "periodic\_z0"}}]\vspace{5pt}
%      \item[{\BCinptext[\texttt{6}]{Set\_BC\_Value = 81}}]
%      \item[{\BCinptext[\texttt{6}]{Periodic\%Matching\_BC\_Name = "periodic\_zp013"}}]
%      \item[{\BCinptext[\texttt{6}]{Periodic\%Matching\_BC\_Idx = 5}}]
%      \item[{\BCinptext[\texttt{6}]{Periodic\%Translation(3) = 0.052}}]
%      %
%    \end{InputExample}
%  \item Taylor-Couette flow\cite{Wang_2021} with streamwise periodicity
%    \begin{InputExample}[before=\footnotesize]
%      %
%      \item[{\BCinptext[\texttt{3}]{Name = "Inflow"}}]
%      \item[{\BCinptext[\texttt{3}]{Set\_BC\_Value = 81}}]
%      \item[{\BCinptext[\texttt{3}]{Periodic\%Matching\_BC\_Name = "Outflow"}}]
%      \item[{\BCinptext[\texttt{3}]{Periodic\%Matching\_BC\_Idx = 4}}]
%      \item[{\BCinptext[\texttt{3}]{Periodic\%Translation(1) = 6.283185307}}]
%      %
%      \item[{\BCinptext[\texttt{4}]{Name = "Outflow"}}]\vspace{5pt}
%      \item[{\BCinptext[\texttt{4}]{Set\_BC\_Value = 81}}]
%      \item[{\BCinptext[\texttt{4}]{Periodic\%Matching\_BC\_Name = "Inflow"}}]
%      \item[{\BCinptext[\texttt{4}]{Periodic\%Matching\_BC\_Idx = 3}}]
%      \item[{\BCinptext[\texttt{4}]{Periodic\%Translation(1) = -6.283185307}}]
%      %
%    \end{InputExample}
%\end{itemize}

%###############################################################################

\IDEA{Each example should include a schematic of the labeled boundary
surfaces to help visualize the periodicity.}

%###############################################################################

